The simulation remains partial. It carries \ValShare{} of counted intersection traffic and under
half at the four busiest junctions. Residential streets and lanes are absent, so any traffic that
really uses them is forced onto modeled roads, and trips to and from the outer valley enter only at
the edge of the box. The demand is from 2011 while the network is recent, the peak hour factor of
0.10 and the give way rule are assumptions, and Ring Road speed is still \ValRingGap~km/h below the
survey. Relative comparisons between interventions are better supported than absolute travel
times.

Route choice is within day dynamic user optimum, rerouting every 600 seconds; a day to day
equilibrium did not converge in the first version and was not retried. Pricing holds demand fixed,
so it can only reroute trips, and no local price elasticity is available. Several analyses from the
first version were not rerun on the new baseline and are omitted: a sweep over assumed lane
capacity, a sensitivity sweep on fitted demand factors, and an elastic demand pricing test.

Deployability differs sharply between levers. No well documented case was found of congestion
pricing implemented and measured in a lower or middle income city with institutional capacity
comparable to Kathmandu's. The closest precedent is Jakarta, where electronic road pricing did not
launch after about a decade, and the occupancy rule used instead was defeated by drivers hiring
passengers \citep{jakartapost2013}. Adaptive signals need no enforcement of driver behavior, and in
this model they deliver most of the combined lever's benefit on their own.

An independent study agrees in direction. \citet{tiwari2024satdobato}, calibrated to field data at
Satdobato and Gwarko with the same GEH statistic, model average delays falling from 226.5 and 201.1
seconds to 51.8 and 52.5 seconds after signal optimization. They also found coordinating the two
intersections worse than optimizing each alone, a reminder that combining measures does not
guarantee more than their sum.
